% !TITLE 正气歌
% !AUTHOR 文天祥
% !DYNASTY 两宋
% !GENRE 五言古诗
% !ORDER 19

\begin{poem}
天地有正气\textsuperscript{1}，杂然赋流形\textsuperscript{2}。
下则为河岳，上则为日星。
于人曰浩然\textsuperscript{3}，沛乎\textsuperscript{4}塞苍冥\textsuperscript{5}。
皇路当清夷\textsuperscript{6}，含和吐明庭\textsuperscript{7}。
时穷节乃见\textsuperscript{8}，一一垂丹青\textsuperscript{9}。

在齐太史简\textsuperscript{10}，在晋董狐笔\textsuperscript{11}。
在秦张良椎\textsuperscript{12}，在汉苏武节\textsuperscript{13}。
为严将军头\textsuperscript{14}，为嵇侍中血\textsuperscript{15}。
为张睢阳齿\textsuperscript{16}，为颜常山舌\textsuperscript{17}。
或为辽东帽\textsuperscript{18}，清操厉冰雪。
或为出师表\textsuperscript{19}，鬼神泣壮烈。
或为渡江楫\textsuperscript{20}，慷慨吞胡羯。
或为击贼笏\textsuperscript{21}，逆竖头破裂。

是气所磅礴\textsuperscript{22}，凛烈万古存\textsuperscript{23}。
当其贯日月，生死安足论。
地维\textsuperscript{24}赖以立，天柱\textsuperscript{25}赖以尊。
三纲\textsuperscript{26}实系命，道义为之根。

嗟予遘阳九\textsuperscript{27}，隶也实不力\textsuperscript{28}。
楚囚缨其冠\textsuperscript{29}，传车\textsuperscript{30}送穷北。
鼎镬甘如饴\textsuperscript{31}，求之不可得。
阴房阒鬼火\textsuperscript{32}，春院閟天黑\textsuperscript{33}。
牛骥同一皂\textsuperscript{34}，鸡栖凤凰食\textsuperscript{35}。
一朝蒙雾露\textsuperscript{36}，分作沟中瘠\textsuperscript{37}。
如此再寒暑，百沴自辟易\textsuperscript{38}。
哀哉沮洳场\textsuperscript{39}，为我安乐国。
岂有他缪巧\textsuperscript{40}，阴阳不能贼\textsuperscript{41}。
顾此耿耿存\textsuperscript{42}，仰视浮云白。
悠悠我心悲，苍天曷有极\textsuperscript{43}。
哲人日已远\textsuperscript{44}，典刑在夙昔\textsuperscript{45}。
风檐展书读\textsuperscript{46}，古道照颜色\textsuperscript{47}。
\end{poem}

\begin{annotations}
\notetext{1}{正气：刚正宏大的气节。文天祥认为天地间存在一种正气，赋予万物以形态——在地为山河，在天为日月，在人为浩然之气。}
\notetext{2}{杂然赋流形：纷纷地赋予万物各种形态。杂然，纷繁多样；赋，赋予；流形，流动变化的形体。}
\notetext{3}{浩然：浩然之气，语出《孟子》"我善养吾浩然之气"，指一种充塞天地的至大至刚之气。}
\notetext{4}{沛乎：充沛盛大的样子。}
\notetext{5}{苍冥：苍天、天空。冥（míng），深远幽暗，指高远的天空。}
\notetext{6}{皇路当清夷：国家的政治道路正应当清明平坦。皇路，国运；清夷，清平。}
\notetext{7}{含和吐明庭：含着和气、吐露于圣明的朝廷。有气节的人在太平年代能为朝廷贡献才智。}
\notetext{8}{时穷节乃见：时局到了穷途末路的时候，人的气节才会显露出来。见（xiàn），显现。这是全文的纲领。}
\notetext{9}{丹青：红色和青色颜料，代指画像、史册。这些有节操的人被绘成肖像流传后世。}
\notetext{10}{齐太史简：春秋时，齐国太史在竹简上直书"崔杼弑其君"，崔杼连杀三位太史，第四位仍坚持如实记录。简，竹简。}
\notetext{11}{董狐笔：春秋时晋国史官董狐，直书"赵盾弑其君"，被孔子称赞为"书法不隐"的良史。}
\notetext{12}{张良椎：张良为报韩国被秦所灭之仇，在博浪沙用大铁椎狙击秦始皇，未中。椎（chuí），铁锤。}
\notetext{13}{苏武节：苏武出使匈奴被扣留十九年，手持汉朝使节（符节）不放，节上的旄毛都掉光了，始终不降。节，使臣的凭证。}
\notetext{14}{严将军头：三国时严颜被张飞俘虏，说"但有断头将军，无降将军"。}
\notetext{15}{嵇侍中血：西晋嵇绍在战乱中以身护卫晋惠帝，被乱军杀害，鲜血溅在惠帝龙袍上。事后惠帝不让洗袍，说"此嵇侍中血，勿洗也"。}
\notetext{16}{张睢阳齿：唐代安史之乱中，张巡死守睢阳，城破被俘，大骂叛军，牙齿都被敲碎。睢（suī）阳。}
\notetext{17}{颜常山舌：唐代颜杲卿任常山太守，安史之乱中城破被俘，骂不绝口，叛军割了他的舌头。}
\notetext{18}{辽东帽：东汉末年管宁避乱辽东，常戴皂帽、穿布衣，终身不仕魏。"清操厉冰雪"——清白的操守比冰雪还洁净。}
\notetext{19}{出师表：诸葛亮北伐前写给后主刘禅的《出师表》，情辞恳切，"鬼神泣壮烈"——鬼神读了都为之落泪。}
\notetext{20}{渡江楫：东晋祖逖率军渡江北伐，在长江中敲打船桨发誓收复中原。楫（jí），船桨。}
\notetext{21}{击贼笏：唐代段秀实用笏板猛击叛臣朱泚的额头，朱泚头破血流。笏（hù），大臣上朝拿的板子；逆竖，对叛贼的蔑称。}
\notetext{22}{磅礴：气势盛大、充满。正气盛大到极点。}
\notetext{23}{凛烈万古存：庄严猛烈地万古长存。凛（lǐn），庄严不可侵犯。}
\notetext{24}{地维：维系大地的绳索。古人认为大地由四角的绳子拉着，才能稳固。}
\notetext{25}{天柱：支撑天空的柱子。古人认为天由柱子撑着。正气是天地立命的根本。}
\notetext{26}{三纲：君为臣纲、父为子纲、夫为妻纲。三纲靠道义维系，道义的根本就是正气。}
\notetext{27}{嗟予遘阳九：可叹我碰上了厄运。嗟（jiē），叹息；遘（gòu），遭遇；阳九，古代术数认为的大灾难年数。指宋朝灭亡。}
\notetext{28}{隶也实不力：我这个做臣子的实在没有能力挽回局面。文天祥自责。}
\notetext{29}{楚囚缨其冠：像楚囚一样还戴着南冠。楚国俘虏钟仪被囚时仍戴楚国帽子。文天祥被俘后仍保持宋朝衣冠。缨，帽带。}
\notetext{30}{传车：驿站的囚车。被用囚车押送到极北的苦寒之地（指元大都，今北京）。}
\notetext{31}{鼎镬甘如饴：大锅烹煮的酷刑像糖浆一样甜。鼎镬（huò），古代烹人的刑具；饴（yí），糖浆。意为宁受极刑也不屈服。}
\notetext{32}{阴房阒鬼火：阴暗的牢房里寂静无人，只有磷火幽幽地飘。阒（qù），寂静。}
\notetext{33}{春院閟天黑：春天的院子被锁得黑暗不透光。閟（bì），关闭。}
\notetext{34}{牛骥同一皂：老牛和千里马关在同一个槽里。骥（jì），千里马；皂（zào），马槽。比喻贤愚不分。}
\notetext{35}{鸡栖凤凰食：凤凰在鸡窝里觅食。比喻高洁之人落在卑贱环境中。}
\notetext{36}{一朝蒙雾露：有一天染上瘴气疫病。蒙雾露，被雾气露水所伤。}
\notetext{37}{分作沟中瘠：按理说会死在沟壑中变成枯骨。分（fèn），按理；瘠（jí），枯骨。}
\notetext{38}{百沴自辟易：各种邪气（瘟疫、瘴气）都自动退避了。沴（lì），灾气、恶气；辟易，退避。正气护身，百毒不侵。}
\notetext{39}{沮洳场：低湿泥泞的地方，指阴暗潮湿的牢房。沮洳（jù rù），低湿。}
\notetext{40}{缪巧：机巧、窍门。缪（miù），巧诈。并没有什么奇特的技巧能活下来。}
\notetext{41}{阴阳不能贼：阴阳二气（冷热、疾病）也不能伤害我。贼，伤害。正气在体内护佑。}
\notetext{42}{耿耿存：一颗耿耿忠心犹在。耿耿，明亮的样子，形容心光明磊落。}
\notetext{43}{苍天曷有极：苍天啊，这悲伤什么时候才有尽头？曷（hé），何。典出《诗经》"悠悠苍天，曷其有极"。}
\notetext{44}{哲人日已远：古代贤哲们一天比一天远了。}
\notetext{45}{典刑在夙昔：但他们的榜样还在从前的事迹中。典刑，典范、楷模；夙昔，往日。}
\notetext{46}{风檐展书读：在风吹的屋檐下展开书卷来读圣贤的事迹。}
\notetext{47}{古道照颜色：古人的道义之光映照在我的脸上。道义的光芒穿越时空，照亮了一个临死者的面容。}
\end{annotations}
